\SourceLeaf{9}
\Article{2}
\par
Considerabimus generalissime corpus finitum figurae cuiuscunque, a reliquo spatio infinito per superficiem unam continuam vel plures continuas interque se discretas separatum (si forte corpus cavitatem unam pluresve includat), quarum complexum simpliciter superficiem corporis dicemus. Concipiatur haec superficies in infinita elementa $\mathrm{d}s$ divisa; sit $P$ punctum elementi $\mathrm{d}s$, cuius coordinatae ad tria plana inter se perpendicularia relatae denotentur per $x,y,z$. Sint $PX,PY,PZ$ rectae axibus coordinatarum resp. parallelae, atque in plagas eas directae, versus quas coordinatae incrementa positiva capere supponuntur, porro sit $PQ$ superficiei normalis extrorsumque directa. Sit $M$ punctum attractum ubicunque libet situm, ipsius coordinatae $a,b,c$, atque distantia $PM$ (semper positive accipienda) $=r$. Angulos quos facit recta $PM$ cum $PX,PY,PZ$ denotabimus per $MX,MY,MZ$, angulosque inter $PQ$ atque $PX,PY,PZ,PM$ per $QX,QY,QZ,QM$. Haec omnia ad puncta superficiei indefinite referuntur: quoties de pluribus punctis superficiei determinatis agendum erit, iisdem characteribus accentibus distinctis utemur.

\Article{3}
\par
Concipiatur planum axi coordinatarum $x$ normale, ita tamen, ut si ipsius aequatio exhibeatur per $x=\alpha$, $\alpha$ sit minor quam valor minimus coordinatae $x$ in superficie corporis. Corpus in hoc planum proiectum figuram finitam ibi designabit, quam in elementa infinita $\mathrm{d}\Sigma$ dispertitam supponemus. In elementi $\mathrm{d}\Sigma$ puncto $\Pi$ erigatur perpendiculum (sive axi coordinatarum $x$ parallelum), quod secet corpus in punctis $P',P'',P'''$ etc.: horum punctorum multitudo manifesto erit par. Erigantur etiam perpendicula ad planum in singulis punctis circumferentiae elementi $\mathrm{d}\Sigma$, quae formabunt superficiem cylindricam sensu latiori, atque e superficie corporis elementa $\mathrm{d}s',\mathrm{d}s'',\mathrm{d}s'''$ etc. rescindent. Elementum $\mathrm{d}\Sigma$ erit proiectio singulorum elementorum $\mathrm{d}s',\mathrm{d}s'',\mathrm{d}s'''$ etc., unde patet esse $\mathrm{d}\Sigma=\pm\mathrm{d}s'.\cos QX'=\pm\mathrm{d}s''.\cos QX''=\pm\mathrm{d}s'''.\cos QX'''$ etc., signo superiori vel inferiori valente, prout cosinus anguli acuti vel obtusi adest. Quoniam vero manifesto perpendiculum in $P'$ corpus ingreditur, in $P''$ e corpore exit, in $P'''$ rursus intrat etc., facile perspicitur, $QX'$ obtusum esse, $QX''$ acutum, $QX'''$ obtusum etc., ita ut habeatur
